% Notes of the Montreal seminar, 1970 (cote 129, typed part): archive cote 129, site page 64 (PDF page 65).
% typescript. Draft compared: none.
% AI-transcribed by Claude Opus 5.5 (claude-opus-5-5, adaptive thinking, effort high), Message Batches API, batch mtl129_full (msgbatch_01AGnDKHwy6ATCM7ehUa6263), 2026-11-13; not yet checked by eye.
\scanpage{129/64}
\mspage{\hand{III-9}}

comme $G'$ est p-divisible, $G(1)$ est une extension de groupes finis
localement libres et est donc de même nature.

b) Un quotient d'un groupe p-divisible (resp. de p-torsion) est
p-divisible (resp. de p-torsion). Comme précédemment, $G''(1)$ est
quotient de deux schémas en groupes finis localement libre, ce qui im-%
plique que $G''(1)$ l'est aussi (II, 1.6).

\uline{Remarque}. Si $G$ et $G''$ sont des groupes de BT, il est en général
faux que le noyau soit un groupe de BT. Par exemple, on a la suite
exacte $0 \to G(1) \to G \xrightarrow{p} G \to 0$ où $G(1)$ n'est pas p-divisible!

5.3. \uline{Rang (ou hauteur) d'un groupe de Barsotti-Tate}. $G(1)$ est un
groupe fini localement libre annulé par $p$. Pour tout $s \in S$, on a
donc
\[ \mathrm{rang}\ (G(1)_{s}) = p^{r(s)} \]
où la fonction $s \mapsto r(s)$ est localement constante. Des suites exactes
$0 \to G(1) \to G(n) \xrightarrow{p} G(n-1) \to 0$ on tire immédiatement que
\[ \mathrm{rang}\ (G(n)_{s}) = [\mathrm{rang}\ (G(1)_{s})]^{n} = p^{nr(s)} \]

L'entier $r(s)$ décrit donc le rang de tous les $G(n)$. Par
définition $r(s)$ s'appelle \uline{le rang} \uline{(ou la hauteur)} \uline{du groupe de Barsotti-}
\uline{Tate} $G$ \uline{au point} $s \in S$.

5.4. On a décrit la catégorie des groupes de BT sur $S$ à l'aide
de systèmes co-p-adiques. Signalons qu'on peut le faire à l'aide de
\uline{systèmes p-adiques:} Si $G$ est un groupe de Barsotti-Tate, on lui fait

\margin{64}
